\SourceLeaf{532}
\fancyhead[L]{\small CONNAISSANCE DES TEMS.}
\begin{center}\small Göttingische gelehrte Anzeigen. Stück 153. pp. 1526..1528. September 24, 1810.\\[3pt]\rule{.72\linewidth}{.3pt}\end{center}

\emph{Connaissance des tems ou des mouvemens célestes, à l'usage des astronomes et des navigateurs, pour l'an} 1811, \emph{publiée par le bureau des longitudes.} 1809. Octav. Paris.

We pass over the calendar, whose arrangement has remained unchanged, and report only the contents of the supplements. Three quarters of them consist of the astronomical observations made at the Imperial Observatory in Paris in 1807 and 1808, whose importance we hardly need to point out here. In addition to the Sun, the Moon and the principal fixed stars, all the older planets were observed very diligently. Among the new ones there are also numerous observations from 1807, especially of Vesta and Ceres. Pallas was observed several times, Juno not at all. In 1808 we find observations of only Vesta and Juno among the new planets, and of the latter only once. We are struck by the sometimes very considerable irregularities in the rate of the pendulum clock recorded in this journal. They tried to remedy them by shifting the pendulum bob more frequently. Vesta and the great comet of 1807 were observed with the parallactic instrument. \textsc{Bouvard} gives his elements for the comet. Next comes a paper by \textsc{Laplace} on the decrease in the obliquity of the ecliptic, based on old observations. Particularly interesting here is the discussion of some ancient Chinese observations, the oldest of which was made 1100 years before our era. Like the Greek, Arabic and Persian determinations, they all agree with the decrease that \textsc{Laplace} determined from theory about as well as could be expected. Another paper by \textsc{Laplace} concerns the rotation of Saturn’s ring, on which, as is well known, \textsc{Schröter}’s observations gave a result entirely different from \textsc{Herschel}’s and incompatible with theory. In the former observations there appeared to be irregularities on the ring that did not change position. \textsc{Laplace} explains these observations by the ingenious hypothesis that the ring actually consists of many concentric rings lying in different planes, which could indeed display such irregularities. So long as the planes of the individual rings did not change their positions appreciably, these irregularities would of course appear unchanged even though the rings were rotating. What remains unexplained is why the two ansae did not present the same appearance. We may hope that repeating these delicate observations in the future will clarify this point. — A paper that follows, by \textsc{Delambre}, on determining latitude and time from the observed altitudes of two known stars, is intended chiefly as a commentary on a programme by \textsc{Gauss} that was also announced in our journal (see Gött. gel. Anz. 1808, p. 1945 [p. 37 of this volume]). Yet \textsc{Delambre} seems to have known this short work only from the abstract in the \emph{Monatliche Correspondenz} and therefore does not judge its purpose quite correctly. \textsc{Delambre} also cites the solution to the same problem that \textsc{Mollweide} communicated in the \emph{Monatl. Correspondenz}. Although on other grounds he ranks it below the trigonometric solution, he praises the ease with which it allows one to choose the correct one of the two roots of the equation. We cannot agree with this judgment. The criterion given by \textsc{Mollweide} is indeed simple, but inadmissible and foreign to the problem. — Calculation by \textsc{van Beek Calkoen} of the observations of the solar eclipse of 16 June 1806 at Utrecht, Milan, Munich and Lilienthal. — \textsc{Burckhardt} on the various ways astronomers have observed the Sun. — Three more short papers by the same astronomer on the comets of 1701 and 1772 and on the \SourceLeaf{533}gridiron pendulum. — Among works relating to astronomy, the only one announced this time is \textsc{d’Entrecasteaux}’s voyage. — An extract from the meteorological journal kept at the observatory by \textsc{Bouvard} in 1807 closes the volume.
